\section*{Abstract}

The CHESS (CubeSat for High-Energy Spectroscopy) mission, currently
under development at the EPFL Spacecraft Team, requires an Electrical
Power System (EPS) capable of managing solar-array energy capture,
battery charging, load distribution, and autonomous safety
behaviour. This report describes the design, implementation, and
hardware-in-the-loop bench-test plan for the firmware that runs on
the EPS microcontroller --- a Microchip SAMD21 --- and decides, in
real time, how the satellite's electrical subsystem operates.

The firmware is bare-metal C running in a single super-loop without
an operating system or vendor hardware abstraction layer. It
implements a Power Conditioning Unit (PCU) state machine with four
normal operating modes and a safe-mode override, an Incremental
Conductance Maximum Power Point Tracking (MPPT) algorithm, a
complementary 300\,kHz PWM driver for the EPC2152 GaN half-bridge
buck converter, the CHIPS framing protocol that the satellite uses
for all internal serial communication, and a sensor abstraction that
lets the same code consume either injected test values or real ADC
and I\textsuperscript{2}C readings.

A hardware-in-the-loop test harness, built from a Python web
application, an ESP32 serial bridge, and the real PCU testing board
V4.1, is presented. Twelve preset scenarios exercise every base case
of the state machine; an MPPT convergence page demonstrates the
algorithm tracking a configurable solar-panel current--voltage curve
in closed loop; a manual-control page lets an operator drive every
board output independently. A nine-test bench-measurement program is
described that produces the physical evidence required to support
the claim that each firmware capability was verified on the real
hardware.

The report closes by enumerating the differences between the present
simulator-grade implementation and the flight-grade firmware the
mission ultimately requires, and outlines a clearly-staged path
between the two.
